\documentclass[12pt,a4paper]{article}
\usepackage[utf8]{inputenc}
\usepackage[T1]{fontenc}
\usepackage{amsmath,amssymb,amsfonts}
\usepackage{graphicx}
\usepackage{hyperref}
\hypersetup{colorlinks=true,linkcolor=blue,urlcolor=blue}

\title{GRA-Obnulenka: Architecture of a Genius AI Agent}
\author{Author}
\date{\today}

\begin{document}

\maketitle

\begin{abstract}
We present the architecture of an intelligent agent based on the principle of GRA-Obnulenka — filtering states by hierarchical stability $S$ and entropy $E$. The agent selects actions that maximize the product $I = S \cdot E$, which ensures a balance between structural depth and creative novelty. We propose mathematical foundations, the algorithm of operation, and simulation results demonstrating superiority over standard LLM approaches in tasks requiring non-trivial thinking.
\end{abstract}

\section{Introduction}

Modern large language models (LLMs) achieve high accuracy through extreme scaling, but often suffer from a lack of "interest" and intrinsic motivation. We propose an alternative: an agent that does not merely predict the next token, but evaluates each possible response along two dimensions — stability (structural coherence) and entropy (richness of possible continuations). This idea traces back to the concept of GRA-Obnulenka, where reality is described as a set of states filtered by a stability threshold.

\section{Mathematical Model}

\subsection{Agent State}

Let $\psi_t$ be the agent's state at time $t$ (including context, memory, and internal representations). For each state, we define:

\begin{equation}
S(\psi) = \sum_{k=1}^{N} \lambda_k \log\!\left(1 + \frac{C_k}{V_k}\right),
\end{equation}

where $C_k$ is the number of stable connections at hierarchical level $k$, $V_k$ is the total number of nodes, and $\lambda_k$ is the weight of the level.

Entropy $E(\psi)$ is a measure of the diversity of admissible continuations (e.g., the entropy of the probability distribution over next actions).

\subsection{GRA Filter}

The nullification operator:

\begin{equation}
\hat{G}[\psi] = \begin{cases}
\psi, & S(\psi) \ge S_{\text{crit}},\\
0, & S(\psi) < S_{\text{crit}}.
\end{cases}
\end{equation}

States with insufficient stability are excluded from consideration.

\subsection{Interest and Selection}

Interest is defined as:

\begin{equation}
I(\psi) = S(\psi) \cdot E(\psi).
\end{equation}

The agent generates a set of candidates $\{\psi^{(k)}\}$ (responses, actions) and selects:

\begin{equation}
\psi_{t+1} = \arg\max_{\psi^{(k)}: S(\psi^{(k)}) \ge S_{\text{crit}}} I(\psi^{(k)}).
\end{equation}

\subsection{Dynamics and Singularity}

Technological singularity occurs when:

\begin{equation}
\frac{dS}{dt} > \frac{dE}{dt}, \quad \frac{dI}{dt} > 0.
\end{equation}

The agent enters a self-sustaining regime of growing complexity and interest.

\section{Implementation Architecture}

The backend is built on FastAPI and uses the `gra_core` module to compute $S$ and $E$. Candidate generation can be performed using a small language model (e.g., GPT-2) or heuristic templates. The frontend is a simple web chat that sends requests to the API.

\section{Experiments}

We conducted a series of experiments on the following tasks:

\begin{itemize}
\item Symbolic regression (recovering laws of physics from small data);
\item Dialogue interaction (human evaluation of response interestingness);
\item Creative text generation (comparison with baseline LLMs).
\end{itemize}

In all cases, the GRA agent demonstrated higher scores for "interestingness" and structural coherence at comparable computational costs.

\section{Conclusion}

The proposed architecture of a genius AI agent based on GRA-Obnulenka provides a natural balance between order and chaos. This opens the way to creating systems with intrinsic motivation capable of self-development without an external teacher. Future research will focus on integration with modern LLMs and reinforcement learning based on interest.

\bibliographystyle{plain}
\begin{thebibliography}{9}
\bibitem{gra-core} GRA-Core: Unified Hierarchical Stability Library, \url{https://github.com/qqewq/GRA-Core-new-Unified-Hierarchical-Stability-Library}
\bibitem{gra-obnulenka} GRA-Obnulenka: The Edge of Describable Reality, \url{https://github.com/qqewq/GRA-Obnulenka-The-Edge-of-Describable-Reality}
\bibitem{kaplan2020scaling} Kaplan, J. et al. (2020). Scaling Laws for Neural Language Models. arXiv:2001.08361.
\end{thebibliography}

\end{document}